% Material suplementar (v3-12): Table S1, hiperparametros dos dois modelos treinados.
% As duas metades vem de ../tables_v3-12/ (geradas por scripts/v3_12_gerar_tabelas.py, a partir de
% fase4/B7.2_hiperparametros/b7_2_folha_hiperparametros.json); nada e digitado aqui.
% Compilar com o diretorio atual = suplementar/:
%   /trabalho/ambientes/texlive-2025/bin/x86_64-linux/pdflatex Table_S1.tex
\documentclass[10pt,a4paper]{article}
\usepackage[a4paper,margin=2.2cm]{geometry}
\usepackage[T1]{fontenc}
\usepackage[utf8]{inputenc}
\usepackage{amsmath,amssymb}
\usepackage{booktabs,array}
\usepackage{microtype}
\pagestyle{empty}
\setlength{\parindent}{0pt}
\renewcommand{\arraystretch}{1.12}
\begin{document}

{\large\bfseries Table S1. Training settings of the two models}

\medskip
\small
Graph regressor and perceptron, 60 runs each; ``same'' marks a value identical in the two models.
Architecture and loss (first half; the optimisation, selection, baselines and cost follow on the next page).

\medskip
% gerado por scripts/v3_12_gerar_tabelas.py -- nao editar a mao
% TA1: fase4/B7.2_hiperparametros/b7_2_folha_hiperparametros.json (60 + 60 corridas); 'same' = valor identico nos dois bracos; so o limiar do corte de gradiente (a fracao de passos cortados nao e impressa); colunas p{} (largura do texto); grupos: Architecture, Loss
\begin{tabular}{>{\raggedright\arraybackslash}p{0.22\linewidth}>{\raggedright\arraybackslash}p{0.43\linewidth}>{\raggedright\arraybackslash}p{0.25\linewidth}}
\toprule
Setting & Graph regressor & Perceptron \\
\midrule
\multicolumn{3}{l}{\textit{Architecture}} \\
Layers & 4 (message passing) & 5 (encoder) \\
Neighbourhood & one hop per relation, full neighbourhood (no sampling) & node batches, no graph \\
Operators & transmitter$\to$terrain and terrain$\to$terrain (Moore neighbourhood with self-loop): GATv2, 4 heads concatenated, edge dimension 2; terrain$\to$transmitter: SAGE with mean aggregation, no gradient & none (no edges) \\
Relation aggregation & sum & -- \\
Layer & LeakyReLU (slope 0.2), dropout, residual connection, LayerNorm & linear, BatchNorm, LeakyReLU (slope 0.2), dropout (not after the last block) \\
Attention heads & 4 & -- \\
Width & 256 (hidden) & 654, 654, 600, 636, 256 (encoder) \\
Input features per terrain node & 18 & same \\
Input features per transmitter & 6 & not used \\
Output channels & 5 & same \\
Dropout & 0.1 & same \\
Parameters, nominal & 1\,812\,515 & 1\,484\,067 \\
Parameters receiving gradient & 1\,484\,037 & same \\
Width-sensitivity arm (10 runs) & -- & widths 712, 712, 720, 688, 256; 1\,812\,515 parameters \\
\midrule
\multicolumn{3}{l}{\textit{Loss}} \\
Prediction term & Huber loss with threshold 5~dB on each of the 5 channels, averaged over the seed nodes (sentinel nodes included) & same \\
Channel weights & uniform, $1/5$ each (fixed) & same \\
Distance-gradient penalty & weight 0.05; $\operatorname{mean}_{(i,j):\,\mathrm{dist}_i<\mathrm{dist}_j}\max\{0,\;\hat P_j-\hat P_i+1~\mathrm{dB}\}$ over $K$ node pairs drawn in each batch, $K$ being the smaller of 512 and half the batch size rounded down, $\mathrm{dist}$ being the distance to the nearest transmitter & same \\
Variance penalty & weight 0.02; $\max\{0,\;\operatorname{sd}(P)-\operatorname{sd}(\hat P)\}^{2}$ over the batch, target without gradient & same \\
Shadowing$\times$NDVI penalty & weight 0.03; $\max\{0,\;0.15-\operatorname{corr}(\hat L_{\mathrm{veg}},\mathrm{NDVI})\}$ over the batch, NDVI without gradient, $\hat L_{\mathrm{veg}}$ being the vegetation component of the decoded path loss; identically zero in these runs & same \\
Free-space loss penalty & weight 0.05; $\operatorname{mean}\max\{0,\;\mathrm{FSPL}_{900}(\mathrm{dist})-\hat L_{\mathrm{tot}}\}$ over the batch, $\hat L_{\mathrm{tot}}$ being the decoded total path loss & same \\
\bottomrule
\end{tabular}


\clearpage
{\large\bfseries Table S1 (continued)}

\medskip
\small
% gerado por scripts/v3_12_gerar_tabelas.py -- nao editar a mao
% TA1: fase4/B7.2_hiperparametros/b7_2_folha_hiperparametros.json (60 + 60 corridas); 'same' = valor identico nos dois bracos; so o limiar do corte de gradiente (a fracao de passos cortados nao e impressa); colunas p{} (largura do texto); grupos: Optimisation, Selection, Baselines, Hardware and cost
\begin{tabular}{>{\raggedright\arraybackslash}p{0.22\linewidth}>{\raggedright\arraybackslash}p{0.43\linewidth}>{\raggedright\arraybackslash}p{0.25\linewidth}}
\toprule
Setting & Graph regressor & Perceptron \\
\midrule
\multicolumn{3}{l}{\textit{Optimisation}} \\
Optimiser & AdamW, weight decay $10^{-5}$ & same \\
Initial learning rate & $10^{-3}$ & same \\
Learning-rate schedule & cosine annealing with warm restarts, $T_0=10$ epochs, $T_{\mathrm{mult}}=2$, minimum rate $10^{-6}$, stepped per epoch; no restart falls within the 8 epochs & same \\
Gradient clipping, norm \mbox{threshold} & 0.5 & same \\
Mixed precision & automatic mixed precision (CUDA) & same \\
Batch size (nodes) & 12288 & same \\
Epochs & 8 & same \\
Early stopping & none & same \\
Gradient steps per run, median (range) & 5912 (5880--5920) & same \\
\midrule
\multicolumn{3}{l}{\textit{Selection}} \\
Epoch retained & lowest validation MAE of received power & same \\
Test set used in selection & no & same \\
Best epoch, median (range) & 6 (2--8) & 6 (1--8) \\
\midrule
\multicolumn{3}{l}{\textit{Baselines}} \\
Frequency (MHz) & 900 & same \\
Heights (m) & Hata rural and COST-231: $h_{\mathrm{tx}}=30$, $h_{\mathrm{rx}}=1.5$ & same \\
COST-231 variant & suburban ($C_m=0$) & same \\
Distance & to the nearest transmitter, floored at 1~m & same \\
Calibration & offset equal to the training median of received power plus modelled path loss, over (a) all training nodes, (b) valid training nodes only & same \\
\midrule
\multicolumn{3}{l}{\textit{Hardware and cost}} \\
Hardware and software & NVIDIA GeForce RTX 5070 Ti; PyTorch 2.10.0+cu128; CUDA 12.8 & same \\
Time per run (s), median (range) & 783.05 (775.9--1141.3) & 84.55 (82.5--192.2) \\
Peak GPU memory (MB), median (range) & 4159.9 (4140.7--4187.6) & 316.3 \\
\bottomrule
\end{tabular}


\end{document}
